% main.tex - showcase deck for the beamer-up derivative
% Current checkout note: this demo now exercises the supported \usetheme{UP} interface.
% Direct \usetheme{UU} loading is retained only as deprecated compatibility during the migration period.
% Attribution: derived from the original beamer-uu work by A.J.H. (Jos) Zuijderwijk; independently maintained here by Zildjian E. California.
% Branding note: distributed UP-branded assets should follow the UP branding guidelines and are intended for UP constituents/internal use. No official endorsement is implied.
% This file documents the supported public naming contract for the current migration milestone.
\documentclass[aspectratio=169]{beamer}
\usetheme[showlogo]{UP}

% -- Metadata -----------------------------------
\title[UP Beamer Theme]{Beamer Theme for UP {\footnotesize \quad \textcolor{UPGold}{v1.0.0}}}
\subtitle{A \LaTeX\ presentation template for the University of the Philippines System}
\author[H. Doofenshmirtz]{Dr. Heinz Doofenshmirtz}
\institute{Department of Information and Computing Sciences}
\date{\today}


\begin{document}
% ===========================================
%  TITLE SLIDES  (single-author, then multi-author)
% ===========================================

% -- Single-author title --
\begin{frame}[plain]
  \titlepage
\end{frame}

% -- Multi-author / conference title --
\author[Alice et al.]{%
  Alice\inst{1} \and
  Bob\inst{1} \and
  Eve\inst{2}%
}
\addinstitute{1}{Institute of Computer Science, University of the Philippines Los Banos}
\addinstitute{2}{Department of Computer Science, University of the Philippines Diliman}
\multiauthor
\venue{National Research Conference 2026}
\begin{frame}[plain]
  \titlepage
\end{frame}

% -- Reset to single-author for the rest of the demo --
\author[H. Doofenshmirtz]{Dr. Heinz Doofenshmirtz}
\institute{Department of Information and Computing Sciences}
\clearinstitutes
\singleauthor

% ===========================================
%  TABLE OF CONTENTS
% ===========================================
\begin{frame}{Outline}
  \tableofcontents
\end{frame}

% ===========================================
\section{Getting Started}
% ===========================================

\begin{frame}{Using the Theme}
  Load the theme in your preamble and set your metadata:

  \begin{itemize}
    \item \texttt{\textbackslash usetheme\{UP\}} --- basic usage
    \item \texttt{\textbackslash usetheme[showlogo]\{UP\}} --- logo on every slide
    \item \texttt{\textbackslash usetheme[nl]\{UP\}} --- legacy Dutch-logo compatibility path
  \end{itemize}

  \vskip 6pt
  For conference-style multi-author title pages, use the standard
  \texttt{\textbackslash author} / \texttt{\textbackslash institute}
  with \texttt{\textbackslash and} and \texttt{\textbackslash inst\{n\}},
  then activate with:
  \begin{itemize}
    \item \texttt{\textbackslash multiauthor}
    \item \texttt{\textbackslash venue\{National Research Conference 2026\}} (optional)
  \end{itemize}

  \vskip 6pt
  Legacy note: \texttt{\textbackslash usetheme\{UU\}} is now deprecated compatibility only.
\end{frame}

\begin{frame}{Theme Options at a Glance}
  \begin{columns}[T]
    \begin{column}{0.48\textwidth}
      \begin{block}{Flags} \footnotesize
        Pass as options to \texttt{\textbackslash usetheme[...]}\texttt{\{UP\}}:
        \texttt{nl}, \texttt{showlogo}, \texttt{nosectionpage}.
      \end{block}
      \vskip 6pt
      \begin{block}{Special Frames} \footnotesize
        \texttt{\textbackslash standoutframe\{...\}} for emphasis.\\
        \texttt{\textbackslash thankframe\{...\}\{...\}} to close.\\
        \texttt{\textbackslash sectionframe\{...\}} manual section divider.
      \end{block}
    \end{column}
    \begin{column}{0.48\textwidth} \footnotesize
      \begin{block}{Conference Title}
        \texttt{\textbackslash multiauthor} activates multi-author layout.
        Use \texttt{\textbackslash addinstitute\{n\}\{Name\}} to list affiliations
        (numbered automatically). Call \texttt{\textbackslash clearinstitutes} to reset.
      \end{block}
    \end{column}
  \end{columns}

  \vskip 10pt
  The \texttt{showlogo} flag is active in this demo, check out the top right corner of this slide.
\end{frame}

% Manual section divider demo (same style as automatic \AtBeginSection pages)
\sectionframe{This is a manual \textbackslash sectionframe}

% ===========================================
\section{Colour Palette}
% ===========================================

\begin{frame}{UP Institutional Colours}
  \begin{columns}[T]
    \begin{column}{0.48\textwidth}
      \begin{alertblock}{Core Palette}
        \textcolor{UPMaroon}{\rule{1em}{1em}}~Maroon \texttt{\#751518}\\[3pt]
        \textcolor{UPForestGreen}{\rule{1em}{1em}}~Forest green \texttt{\#0A4424}\\[3pt]
        \textcolor{UPGold}{\rule{1em}{1em}}~Gold \texttt{\#F5AB29}\\[3pt]
        \textcolor{UPBlack}{\rule{1em}{1em}}~Spot black \texttt{\#231F20}\\[3pt]
        \textcolor{white}{\rule{1em}{1em}}~White \texttt{\#FFFFFF}
      \end{alertblock}
    \end{column}
    \begin{column}{0.48\textwidth}
    \begin{exampleblock}{Usage Rules}
        \begin{itemize}
          \item Maroon carries hero surfaces and alerts.
          \item Forest green supports section structure and examples.
          \item Gold stays an accent for rules, bullets, and emphasis.
          \item Spot black anchors body text and labels on gold.
        \end{itemize}
    \end{exampleblock}
    \end{column}
  \end{columns}

  \vskip 8pt
  The public theme now stays inside the UP four-colour system rather than the inherited multi-hue palette.
\end{frame}

% ===========================================
\section{Typography \& Blocks}
% ===========================================

\begin{frame}{Typography}
  \begin{enumerate}
    \item Optima-directed headings interpreted through a restrained sans-serif treatment
    \item \textbf{Palatino} body text for the core reading surface
    \item The supported \texttt{pdfLaTeX} baseline keeps headings on a Helvetica-style sans family for reproducible builds
  \end{enumerate}

  \vskip 10pt
  \begin{exampleblock}{Tip}
    The supported baseline is still \texttt{pdfLaTeX}; the theme now favors the approved UP hierarchy over the older Merriweather/Open Sans pairing.
  \end{exampleblock}
\end{frame}

\begin{frame}{Lists}
  Itemize now uses a single gold marker family across levels:

  \begin{itemize}
    \item First-level items use \textcolor{UPGold}{gold} squares
    \begin{itemize}
      \item Second-level items stay gold and shrink with depth
      \begin{itemize}
        \item Third-level items keep the same accent hue for consistency
      \end{itemize}
    \end{itemize}
    \item Items can contain \textbf{bold}, \emph{italic}, or \alert{alerted} text
    \item Keep bullets concise and let spacing carry hierarchy
  \end{itemize}

  \vskip 8pt
  Enumerate uses dark labels on gold chips:
  \begin{enumerate}
    \item First item
    \item Second item
    \item Third item
  \end{enumerate}
\end{frame}

\begin{frame}{Mathematics}
  The theme supports mathematical typesetting \cite{tai1994mathematical}:

$$
\int_{a}^{b} f(x)\,dx \approx 
\frac{h}{2}\left[f(x_0) + 2\sum_{i=1}^{n-1} f(x_i) + f(x_n)\right]
$$

  \begin{alertblock}{Note}
    This is a note block, intended to grab the reader's attention.
  \end{alertblock}
\end{frame}

% ===========================================
\section{Diagrams \& Code}
% ===========================================

\begin{frame}{Workflow Diagram}
  \begin{figure}
  \centering
  \begin{tikzpicture}[
    node distance=0.6cm and 1.0cm,
    box/.style={
      rectangle, rounded corners=3pt,
      minimum height=2em, minimum width=4.5em,
      text centered, font=\sffamily\footnotesize,
      draw=none
    },
    arr/.style={->, >=stealth, semithick, UPBlack80},
    scale=0.95, transform shape
  ]
    \node[box, fill=UPGold, text=UPBlack] (data) {Data};
    \node[box, fill=UPForestGreen, text=white, right=of data] (preprocess) {Pre-process};
    \node[box, fill=UPMaroon, text=white, right=of preprocess] (model) {Model};
    \node[box, fill=UPForestGreen!85!UPBlack, text=white, right=of model] (evaluate) {Evaluate};
    \node[box, fill=UPMaroon!85!UPBlack, text=white, right=of evaluate] (deploy) {Deploy};

    \draw[arr] (data) -- (preprocess);
    \draw[arr] (preprocess) -- (model);
    \draw[arr] (model) -- (evaluate);
    \draw[arr] (evaluate) -- (deploy);

    \draw[arr, UPGold!85!UPBlack, dashed]
      (evaluate.south) -- ++(0,-0.7) -| node[below, pos=0.25, font=\tiny\sffamily, UPMaroon] {iterate} (model.south);
  \end{tikzpicture}
  \caption{Example process diagram with a numbered figure caption.}
  \end{figure}

  \vskip 6pt
  TikZ diagrams now stay inside the mapped UP token set.
  Use \texttt{UPGold}, \texttt{UPForestGreen}, \texttt{UPMaroon}, and \texttt{UPBlack} directly.
\end{frame}

\begin{frame}[fragile]{Code Listing}
The \texttt{listings} package is pre-configured with UP colours.
  Specify the language and write code verbatim:

  \begin{lstlisting}[language=Python]
def fibonacci(n):
    """Return the n-th Fibonacci number."""
    if n <= 1:
        return n
    a, b = 0, 1
    for _ in range(2, n + 1):
        a, b = b, a + b  # shift window
    return b

print(fibonacci(10))  # 55
  \end{lstlisting}
\end{frame}

% ===========================================
\section{Tables \& Figures}
% ===========================================

\begin{frame}{Tables}
  \begin{table}
    \centering
    \caption{Theme files overview}
    \begin{tabular}{ll}
      \hline
      \textbf{File} & \textbf{Purpose} \\
      \hline
      \texttt{beamerthemeUP.sty}       & Canonical public loader wrapper \\
      \texttt{beamerthemeUU.sty}       & Legacy implementation entry \\
      \texttt{beamercolorthemeUU.sty}  & Current UP-directed colour theme \\
      \texttt{beamerinnerthemeUU.sty}  & Title page, blocks, special frames \\
      \texttt{beamerouterthemeUU.sty}  & Frame title, footline, logo \\
      \hline
    \end{tabular}
  \end{table}
\end{frame}

\begin{frame}{Figure in a Column Layout}
  \begin{columns}[T]
    \begin{column}{0.46\textwidth}
      Pair figures with text using \texttt{columns}:
      \begin{itemize}
        \item Left column carries the narrative
        \item Right column holds the figure
        \item Caption is numbered automatically
      \end{itemize}

      \vskip 8pt
      Use \texttt{[T]} on \texttt{columns} to top-align both sides,
      and \texttt{\textbackslash centering} inside the figure.
    \end{column}
    \begin{column}{0.50\textwidth}
      \begin{figure}
        \centering
        \begin{tikzpicture}[scale=0.82]
         % bars
          \foreach \x/\h/\col in {
            0.4/1.2/UPMaroon,
            1.1/2.4/UPMaroon,
            1.8/1.8/UPMaroon,
            2.5/2.9/UPForestGreen,
            3.2/2.1/UPForestGreen,
            3.9/3.0/UPForestGreen}{
            \fill[\col!80] (\x,0) rectangle (\x+0.55,\h);
          }
          % axes
          \draw[UPBlack80, thin] (0,0) -- (4.5,0) node[right,font=\tiny\sffamily]{$x$};
          \draw[UPBlack80, thin] (0,0) -- (0,3.2) node[above,font=\tiny\sffamily]{$y$};
          % legend
          \fill[UPMaroon!80]  (0.3,3.05) rectangle (0.65,3.25);
          \node[right,font=\tiny\sffamily] at (0.7,3.15) {Group A};
          \fill[UPForestGreen!80] (2.0,3.05) rectangle (2.35,3.25);
          \node[right,font=\tiny\sffamily] at (2.4,3.15) {Group B};
        \end{tikzpicture}
        \caption{Sample bar chart with UP colours.}
      \end{figure}
    \end{column}
  \end{columns}
\end{frame}

% -- Standout --------------------------
\standoutframe{Sharing science,\\[4pt]\textit{shaping tomorrow}.}

% ===========================================
\section{Closing}
% ===========================================

\begin{frame}{Summary}
  \begin{enumerate}
    \item Supported public loader: \texttt{\textbackslash usetheme\{UP\}}
    \item Palatino-led reading surface with a restrained sans heading hierarchy
    \item Special frames: \texttt{\textbackslash standoutframe}, \texttt{\textbackslash thankframe}, \texttt{\textbackslash sectionframe}
    \item Options: \texttt{nl}, \texttt{showlogo}, \texttt{nosectionpage}
    \item Multi-author slides: \texttt{\textbackslash addinstitute\{n\}\{Name\}} + \texttt{\textbackslash multiauthor}
    \item Legacy \texttt{\textbackslash usetheme\{UU\}} loading remains deprecated compatibility during migration
  \end{enumerate}

  \vskip 12pt
  {\small Source and documentation: \url{https://github.com/zcalifornia-ph/beamer-up}}
\end{frame}

% -- Closing -----------------------------------
\thankframe{Feedback?}{Please contribute \\ \url{https://github.com/zcalifornia-ph/beamer-up}}

% -- references -----------------------------
\begin{frame}[allowframebreaks]{References}
  The Beamer class~\cite{tantau2004} provides the foundation.
  The UP Visual Identity Guidebook~\cite{uubrand2025} defines the colour palette,
  typography, and logo usage rules.

  \vskip 6pt
  \printbibliography
\end{frame}

\end{document}
